% !TEX program = xelatex
% Generated by docs/latex/build.py from the sources pinned in sources.json.
\documentclass[11pt,a4paper]{article}
\usepackage[a4paper,margin=25mm]{geometry}
\usepackage{fontspec}
\setmainfont{lmroman10-regular.otf}[BoldFont=lmroman10-bold.otf,ItalicFont=lmroman10-italic.otf,BoldItalicFont=lmroman10-bolditalic.otf]
\setsansfont{lmsans10-regular.otf}[BoldFont=lmsans10-bold.otf,ItalicFont=lmsans10-oblique.otf]
\setmonofont{JuliaMono-Regular.ttf}[Path=fonts/,Scale=0.83]
\usepackage{amsmath,amssymb,mathtools}
\usepackage{graphicx,calc,longtable,booktabs,array}
\usepackage{fvextra}
\usepackage{xurl}
\usepackage[colorlinks=true,linkcolor=blue,urlcolor=blue,citecolor=blue]{hyperref}
\providecommand{\tightlist}{\setlength{\itemsep}{0pt}\setlength{\parskip}{0pt}}
\setlength{\parindent}{0pt}
\setlength{\parskip}{6pt}
\setlength{\emergencystretch}{6em}
\begin{document}
\section{Angles between lines: written proofs for C1--C5}\label{angles-between-lines-written-proofs-for-c1c5}

For unit representatives of lines through the origin, write \(\theta(x,y)=\arccos|\langle x,y\rangle|\), \(\delta(x,y)=\pi/2-\theta(x,y)=\arcsin|\langle x,y\rangle|\), and \(S=\sum_{i<j}\theta(x_i,x_j)\). Repeated lines are allowed throughout. These are complete written proofs for the five stated cells. C6 is an open general conjecture and is not claimed here. The arguments do not rely on computation.

\subsection{C1: planar lines}\label{c1-planar-lines}

For the \href{https://hackathon.bainsa.ai/p/p1/c1}{C1 statement}, let \(x_1,\ldots,x_N\) represent \(N\ge0\) lines in \(\mathbb R^2\), and put \(S=\sum_{i<j}\arccos|\langle x_i,x_j\rangle|\). The target is

\[
S\le\frac{\pi}{2}\left\lfloor\frac{N^2}{4}\right\rfloor.
\]

Represent each line by a point \(\alpha\) on the circle \(\mathbb R/(\pi\mathbb Z)\), of circumference \(\pi\). Its distance from another point is the shorter circular arc, hence exactly the non-obtuse line angle.

For \(t\) on that circle, let \(A_t\) be the half-open semicircle \([t,t+\pi/2)\), and let \(k(t)\) count the \(N\) line-points in \(A_t\), with multiplicity. Exactly \(k(t)(N-k(t))\) unordered pairs have one point in \(A_t\) and one outside.

For two fixed points at circular distance \(\delta\in[0,\pi/2]\), the set of \(t\) separating them has length \(2\delta\). Indeed the sets of \(t\) for which each point is in \(A_t\) are semicircles whose symmetric difference consists of two intervals of length \(\delta\); this also handles coincident or opposite points. Endpoints have measure zero.

Integrate the pair count over a full period. All functions are finite step functions, so ordinary Riemann integration and finite additivity suffice:

\[
2S=\int_0^\pi k(t)(N-k(t))\,dt
\le \pi\left\lfloor\frac{N^2}{4}\right\rfloor.
\]

The inequality follows from \(4k(N-k)=N^2-(N-2k)^2\) and integrality. Divide by \(2\). Equality is attained by \(\lfloor N/2\rfloor\) copies of one axis and \(\lceil N/2\rceil\) copies of its perpendicular axis.

\subsection{C2: the tridiagonal case}\label{c2-the-tridiagonal-case}

For the \href{https://hackathon.bainsa.ai/p/p1/c2}{C2 statement}, let \(m\ge2\) and \(x_1,\ldots,x_m\) be unit vectors in \(\mathbb R^{m-1}\) with \(\langle x_i,x_j\rangle=0\) whenever \(|i-j|\ge2\). The target is

\[
\sum_{i=1}^{m-1}\arccos|\langle x_i,x_{i+1}\rangle|\le(m-2)\frac{\pi}{2}.
\]

Let \(G\) be the \(m\times m\) Gram matrix. It has diagonal \(1\) and vanishes off the first off-diagonals. Since the vectors lie in \(\mathbb R^{m-1}\), \(G\) is singular. Successive sign changes of the vectors make every adjacent entry \(a_i\) nonnegative; this preserves all absolute inner products, and hence every line angle. Write \(\beta_i=\arcsin(a_i)\in[0,\pi/2]\). The desired inequality is equivalent to \(\sum_i\beta_i\ge\pi/2\).

Suppose instead \(B=\sum_i\beta_i<\pi/2\). Let \(B_i=\beta_1+\cdots+\beta_i\). Set \(b_1=a_1\), and inductively

\[
b_i=\frac{a_i}{\sqrt{1-b_{i-1}^2}}.
\]

We show \(0\le b_i\le\sin B_i<1\), so every denominator exists and is positive. The base case is \(b_1=\sin\beta_1\). At the inductive step, \(b_{i-1}\le\sin B_{i-1}\) implies

\[
\sqrt{1-b_{i-1}^2}\ge\cos B_{i-1}>0.
\]

Since \(B_i<\pi/2\),

\[
\sin B_i\cos B_{i-1}-\sin\beta_i
=\sin B_{i-1}\cos B_i\ge0.
\]

Therefore \(b_i\le\sin\beta_i/\cos B_{i-1}\le\sin B_i<1\).

Now eliminate successive rows/columns of the tridiagonal matrix \(G\). Its \(LDL^{\mathsf T}\) pivots are

\[
d_1=1,\qquad d_{i+1}=1-\frac{a_i^2}{d_i}=1-b_i^2.
\]

All pivots are positive. Explicitly, take \(L\) lower bidiagonal, with diagonal \(1\) and \(L_{i+1,i}=a_i/d_i\). Direct multiplication gives \(G=L\mathrm{diag}(d_1,\ldots,d_m)L^{\mathsf T}\). Thus \(\det G=\prod_i d_i>0\), contradicting singularity. This contradiction proves \(\sum_i\beta_i\ge\pi/2\) and hence C2.

\subsection{\texorpdfstring{C3: \(d+1\) lines in \(\mathbb R^d\)}{C3: d+1 lines in \textbackslash mathbb R\^{}d}}\label{c3-d1-lines-in-mathbb-rd}

For every \(d\ge1\) and \(d+1\) lines, the \href{https://hackathon.bainsa.ai/p/p1/c3}{C3 statement} asks for

\[
S\le\left(\binom{d+1}{2}-1\right)\frac{\pi}{2}.
\]

Equivalently, their total deficiency is at least \(\pi/2\). We prove the following auxiliary bound for every \(m\ge2\): any \(m\) unit vectors whose span has dimension at most \(m-1\) have total deficiency at least \(\pi/2\). Embed their span into \(\mathbb R^{m-1}\); then set \(N=m\) and \(d=m-1\) in the argument below.

\subsubsection{Sparse-maximizer argument}\label{sparse-maximizer-argument}

For fixed \(N,d\), the angle sum attains a maximum on the compact product of \(N\) unit spheres. Among maximizers choose one with the largest number of orthogonal pairs.

For a vector \(x\) with at least one nonorthogonal neighbor, let \(W\) span its orthogonal neighbors. If \(\dim W<d-1\), choose a unit \(y\) orthogonal to \(W\) and to \(x\). Move \(x\) along \(x(t)=x\cos t+y\sin t\). Every existing orthogonality is preserved. For any other vector \(z\) not orthogonal to \(x\), its inner product is \(A\cos t+B\sin t\), with \(A\ne0\). On the interval between its consecutive zeros containing \(t=0\), its line angle is convex. To check this, write the signed inner product as \(r\cos(t-t_0)>0\), \(0<r\le1\), and put \(u=t-t_0\). For \(r<1\),

\[
\frac{d^2}{du^2}\arccos(r\cos u)=\frac{r(1-r^2)\cos u}{(1-r^2\cos^2 u)^{3/2}}\ge0.
\]

For \(r=1\) the angle is an absolute-value function on that interval, also convex. Intersect the closed intervals bounded by the nearest zeros for all nonorthogonal neighbors. The total angle sum is convex there by continuity at the endpoints, and has its global maximum at the interior point \(0\), so it is constant. At an endpoint at least one additional pair becomes orthogonal, while no orthogonal pair is lost. This contradicts the choice of maximizer.

Here \(N>d\), so at least one pair is nonorthogonal. If the configuration spanned dimension \(r<d\), an endpoint \(x\) of such a pair would have \(\dim W\le r-1<d-1\), contradicting the preceding rotation argument. Thus the configuration spans \(\mathbb R^d\). For a vector orthogonal to every other vector, those other vectors consequently span its entire orthogonal complement; its orthogonal-neighbor span has dimension \(d-1\) too. Every vector therefore has at least \(d-1\) linearly independent orthogonal neighbors. Its nonorthogonality graph (an edge means nonzero inner product) has maximum degree at most \(N-d\).

For \(N=d+1\), this graph consists of edges and isolated vertices. The Gram matrix has nullity one. Some two-vertex block must have rank one, hence its two lines coincide and contribute deficiency \(\pi/2\). This proves the auxiliary deficiency bound for an angle maximizer. Every other configuration has angle sum no greater than that maximizer, so its deficiency is no smaller. Taking \(m=d+1\) proves C3. Equality is attained by taking \(d\) mutually orthogonal coordinate axes and repeating one of them: the only positive deficiency is \(\pi/2\). This also covers \(d=1\) and repeated lines.

\subsection{\texorpdfstring{C4: five lines in \(\mathbb R^3\) and six in \(\mathbb R^4\)}{C4: five lines in \textbackslash mathbb R\^{}3 and six in \textbackslash mathbb R\^{}4}}\label{c4-five-lines-in-mathbb-r3-and-six-in-mathbb-r4}

The \href{https://hackathon.bainsa.ai/p/p1/c4}{C4 statement} asks for \(S\le4\pi\) for five lines in \(\mathbb R^3\) and \(S\le13\pi/2\) for six lines in \(\mathbb R^4\). In either case, with \(D=\sum_{i<j}\delta(x_i,x_j)\), this is equivalent to \(D\ge\pi\). We use the sparse-maximizer argument above with \(N=d+2\).

\subsubsection{Sparse graph and small components}\label{sparse-graph-and-small-components}

For \(N=d+2\), the sparse graph consists of paths and cycles. Gram blocks belonging to different components are orthogonal. An irreducible path block has nullity at most one: the first row and the tridiagonal recurrence determine every kernel coordinate from the first coordinate. An irreducible cycle block has nullity at most two: two adjacent coordinates determine all others. The whole matrix has nullity exactly two. If two components are singular, the C3 auxiliary bound gives deficiency at least \(\pi/2\) in each, proving the target. Otherwise one cycle has nullity two, and all other components are nonsingular.

For \(N\le6\) only cycles of lengths \(3,4,5,6\) remain.

A three-cycle of nullity two has rank one: all three lines coincide and its deficiency is \(3\pi/2\). A four-cycle of nullity two has rank two: its opposite pairs are orthonormal bases of the same plane. For every unit planar vector, its acute angles to an orthogonal pair sum to \(\pi/2\). Consequently the four edge deficiencies sum to \(\pi\).

\subsubsection{Two elementary analytic inequalities}\label{two-elementary-analytic-inequalities}

For \(p,q\in[0,\pi/2]\), put

\[
\begin{aligned}h&=\arccos(\cos p\cos q),\\E(p,q)&=\arcsin\!\left(\frac{\sin p\sin q}{1+\cos p\cos q}\right).\end{aligned}
\]

Then

\[
h+E(p,q)\le p+q.\tag{B}
\]

Here is a direct proof, with no geometric theorem required. Set \(A=\cos p\), \(B=\cos q\), \(r=\sin p\sin q\), \(u=AB\), and \(L=\sqrt{1-u^2}\). The identity

\[
(1+u)^2-r^2=(A+B)^2.
\]

gives \(\cos E=(A+B)/(1+u)\). Therefore

\[
(1+u)[\cos(h+E)-\cos(p+q)]=-u(1-A)(1-B)+r(1+u-L)\ge0.
\]

Indeed \(L\le1\) and \(r\ge(1-A)(1-B)\), since \(\sin t\ge1-\cos t\) on \([0,\pi/2]\). Both arguments compared lie in \([0,\pi]\), where cosine is decreasing. This proves (B).

A second inequality, for \(T,p,q\in[0,\pi/2]\), is

\[
\arcsin(\sin p\sin q)\le\arccos(\cos T\cos p)+\arccos(\sin T\cos q)-\frac{\pi}{2}.\tag{C}
\]

Set \(X=\cos T\cos p\), \(Y=\sin T\cos q\), and \(Z=\sin p\sin q\). Direct expansion gives

\[
X^2+Y^2+Z^2+2XYZ-1=-(\cos T\sin p\cos q-\sin T\sin q\cos p)^2\le0.
\]

For nonnegative \(X,Y,Z\le1\) this implies

\[
\sqrt{1-X^2}\sqrt{1-Y^2}\ge Z+XY.
\]

hence \(\cos(\arccos X+\arccos Y)\le-Z\). Since \(\arccos X+\arccos Y\in[0,\pi]\), cosine monotonicity gives their sum at least \(\pi-\arccos Z\). This is (C).

\subsubsection{The five-cycle}\label{the-five-cycle}

Its Gram rank is three. Take \(x_1=e_1,x_3=e_2\), and choose signs of coordinates so that, in absolute value,

\[
x_4=(0,c,\sqrt{1-c^2}),\qquad x_5=(a,0,\sqrt{1-a^2}).
\]

where \(0<a,c<1\). Possible remaining signs do not affect any formula below. Since \(x_2\) is orthogonal to \(x_4\) and \(x_5\), their cross product determines it up to sign. With \(L=\sqrt{a^2+c^2-a^2c^2}\), the absolute edge inner products are

\[
\frac{c\sqrt{1-a^2}}{L},\quad\frac{a\sqrt{1-c^2}}{L},\quad c,\quad\sqrt{1-a^2}\sqrt{1-c^2},\quad a.
\]

Write \(a=\sin p,c=\sin q\). If the first two deficiencies are \(U,V\), then

\[
\begin{aligned}\cos(U+V)&=\frac{\sin p\sin q}{1+\cos p\cos q},\\\sin(U+V)&=\frac{\cos p+\cos q}{1+\cos p\cos q}.\end{aligned}
\]

Thus \(U+V=\pi/2-E(p,q)\). The middle product has deficiency \(\pi/2-\arccos(\cos p\cos q)\). The total edge deficiency is therefore

\[
\pi+p+q-\arccos(\cos p\cos q)-E(p,q)\ge\pi.
\]

by (B). The expressions extend continuously to boundary cases; alternatively those cases break the cycle and were already handled. This completes the five-line \(\mathbb R^3\) bound.

\subsubsection{Removing a vertex from the six-cycle}\label{removing-a-vertex-from-the-six-cycle}

The six-cycle has Gram rank four. Remove \(x_6\) and project \(x_1\) and \(x_5\) orthogonally onto \(x_6^\perp\), normalizing the results. The other three vectors already lie in that hyperplane. The resulting five lines lie in \(\mathbb R^3\).

Only four original deficiencies and three projected deficiencies need comparison. Relabel the local path \((x_2,x_1,x_6,x_5,x_4)\) as \((v_1,v_2,v_3,v_4,v_5)\), and choose signs making its consecutive inner products \(a,b,c,d\) positive. The odd vectors \(v_1,v_3,v_5\) are orthonormal. In an orthonormal basis extending them,

\[
v_2=(a,b,0,u),\qquad v_4=(0,c,d,v),\qquad uv=-bc.
\]

All \(a,b,c,d,|u|,|v|\) are positive for an irreducible six-cycle. Choose the fourth basis direction so \(u>0\) and \(v<0\). Consequently there are \(P,Q,T\in(0,\pi/2)\) with

\[
\begin{aligned}a&=\cos P,&b&=\sin P\cos T,&u&=\sin P\sin T,\\d&=\cos Q,&c&=\sin Q\sin T,&v&=-\sin Q\cos T.\end{aligned}
\]

Let \(B=\arcsin b,C=\arcsin c\). Let \(p,q\) be the acute angles of the normalized projections of \(v_2,v_4\) to \(v_1,v_5\) respectively. Then

\[
\begin{aligned}\cos P&=\cos B\cos p,&\cos Q&=\cos C\cos q,\\\tan B&=\sin p\cot T,&\tan C&=\sin q\tan T.\end{aligned}
\]

The new connecting inner product has absolute value \(\sin p\sin q\). Hence the original local deficiency is \(\pi-P-Q+B+C\), whereas the projected local deficiency is \(\pi-p-q+\arcsin(\sin p\sin q)\).

By (B),

\[
B+p-P\ge E(B,p),\qquad C+q-Q\ge E(C,q).
\]

The tangent relations imply the following exact identities:

\[
\begin{aligned}E(B,p)&=\arccos(\cos T\cos p)-T,\\E(C,q)&=\arccos(\sin T\cos q)-\left(\frac{\pi}{2}-T\right).\end{aligned}
\]

For completeness the first identity follows by substituting \(\sin B\sin T=\cos B\sin p\cos T\) into the cosine addition formula for \(E(B,p)+T\); the result is \(\cos T\cos p\). Both angles are in \([0,\pi]\), making cosine injective. The second is the same computation with \(T\) replaced by \(\pi/2-T\).

Inequality (C) now gives

\[
(B+p-P)+(C+q-Q)\ge\arcsin(\sin p\sin q).
\]

Thus projection does not increase total deficiency. The projected five-line configuration has deficiency at least \(\pi\) by the result just proved. The original six-cycle therefore also has deficiency at least \(\pi\).

Combining all component cases proves both C4 bounds for a sparse angle maximizer. Every other configuration has no larger angle sum, so the bounds hold for all five- and six-line configurations. Equality is attained by coordinate axes with two axes repeated: multiplicities \((2,2,1)\) in \(\mathbb R^3\) and \((2,2,1,1)\) in \(\mathbb R^4\).

\subsection{\texorpdfstring{C5: \(d+2\) lines in \(\mathbb R^d\)}{C5: d+2 lines in \textbackslash mathbb R\^{}d}}\label{c5-d2-lines-in-mathbb-rd}

For every \(d\ge2\) and \(N=d+2\) lines, the \href{https://hackathon.bainsa.ai/p/p1/c5}{C5 statement} asks for

\[
S\le\left(\binom N2-2\right)\frac{\pi}{2}.
\]

Equivalently, \(D=\sum_{i<j}\delta(x_i,x_j)\ge\pi\). The proof extends the preceding six-cycle projection argument to cycles of any length. Its scalar monotonicity step handles a possible strict Cauchy--Schwarz inequality in higher dimension.

\subsubsection{The scalar projection inequality in its general form}\label{the-scalar-projection-inequality-in-its-general-form}

For acute \(B,U\) define

\[
g(B,U)=B+U-\arccos(\cos B\cos U).
\]

For \(B,C,U,V\in(0,\pi/2)\), suppose

\[
\tan B\tan C\le\sin U\sin V.
\]

Then

\[
\arcsin(\tan B\tan C)\le g(B,U)+g(C,V).\tag{1}
\]

The arcsine argument is in \((0,1)\), since \(\sin U\sin V<1\). Put

\[
U_0=\arcsin\!\left(\frac{\tan B\tan C}{\sin V}\right).
\]

The hypothesis and sine monotonicity on \([0,\pi/2]\) give \(0<U_0\le U<\pi/2\). Choose \(T\in(0,\pi/2)\) with

\[
\tan T=\frac{\sin U_0}{\tan B}=\frac{\tan C}{\sin V}.
\]

Thus

\[
\begin{aligned}\sin B\sin T&=\cos B\sin U_0\cos T,\\\sin C\cos T&=\cos C\sin V\sin T.\end{aligned}
\]

The equality-case projection inequality proved in the preceding C4 six-cycle argument gives

\[
\arcsin(\sin U_0\sin V)\le g(B,U_0)+g(C,V).
\]

Its left side is \(\arcsin(\tan B\tan C)\). It remains to compare \(U_0\) with \(U\). For fixed \(B\in(0,\pi/2)\),

\[
\frac{\partial g(B,U)}{\partial U}=1-\frac{\cos B\sin U}{\sqrt{1-\cos^2B\cos^2U}}\ge0.
\]

The denominator is positive, and its square minus \((\cos B\sin U)^2\) equals \(\sin^2 B\). Consequently \(g(B,U_0)\le g(B,U)\), proving (1).

\subsubsection{Projection of a cycle in any dimension}\label{projection-of-a-cycle-in-any-dimension}

Consider an irreducible nonorthogonality cycle of length \(k\ge6\) and rank \(k-2\). Remove one vertex and project its two neighbors onto its orthogonal complement, normalizing both. All remaining vertices were already orthogonal to the removed vertex.

Write the five-vertex path around the removed vertex as \((v_1,v_2,v_3,v_4,v_5)\), with \(v_3\) removed. Orient the representatives along this path so that the four consecutive inner products \(a,b,c,d\) are positive. Because \(k\ge6\), the odd vectors \(v_1,v_3,v_5\) are orthonormal and \(v_2\perp v_4\). Decompose

\[
\begin{aligned}v_2&=av_1+bv_3+u,\\v_4&=cv_3+dv_5+v,\end{aligned}
\]

where \(u,v\) are orthogonal to the span of \(v_1,v_3,v_5\). Then

\[
\begin{aligned}a^2+b^2+\|u\|^2&=1,\\c^2+d^2+\|v\|^2&=1,\\\langle u,v\rangle&=-bc.\end{aligned}
\]

Cauchy--Schwarz yields \(bc\le\|u\|\|v\|\). Since \(a,b,c,d>0\), the residual vectors are nonzero and

\[
A=\sqrt{1-b^2}>0,\qquad Q=\sqrt{1-c^2}>0.
\]

Set

\[
\begin{aligned}B&=\arcsin b,&C&=\arcsin c,\\\cos U&=a/A,&\sin U&=\|u\|/A,\\\cos V&=d/Q,&\sin V&=\|v\|/Q.\end{aligned}
\]

These specify angles \(B,C,U,V\in(0,\pi/2)\). Moreover,

\[
\tan B\tan C=\frac{bc}{AQ}\le\frac{\|u\|\|v\|}{AQ}=\sin U\sin V.
\]

Only four old deficiencies and three new deficiencies change. The original local contribution is

\[
\arcsin a+B+C+\arcsin d=\pi-\arccos(\cos B\cos U)-\arccos(\cos C\cos V)+B+C.
\]

The normalized projections of \(v_2,v_4\) have inner products \(a/A\) with \(v_1\), \(d/Q\) with \(v_5\), and \(-bc/(AQ)\) with each other. Their local contribution is therefore

\[
\pi-U-V+\arcsin(\tan B\tan C).
\]

Inequality (1) says that the new contribution is at most the old contribution. Every other surviving inner product is unchanged: vectors outside the two projected neighbors are orthogonal to \(v_3\), and all nonincident products with either projected neighbor were zero. The removed vertex had nonzero products only with those two neighbors. Thus projection does not increase total deficiency.

The original cycle spans a space of dimension \(k-2\). After projection, the \(k-1\) remaining lines lie in a space of dimension at most \(k-3\). No assertion about their new Gram rank or graph structure is needed.

\subsubsection{Induction and component accounting}\label{induction-and-component-accounting}

We prove the desired bound by strong induction on \(N\ge4\). The sparse-maximizer argument proved above applies in every finite dimension: some angle maximizer has at most \(N-d=2\) nonorthogonal neighbors per line.

Its nonorthogonality graph has only path and cycle components. The Gram matrix is block diagonal across these components. Its rank is at most \(d=N-2\), so its total nullity is at least two; full spanning is unnecessary here. An irreducible path block has nullity at most one; an irreducible cycle block has nullity at most two. The ordered-band kernel recurrences were established in the C4 argument above.

If at least two components are singular, each has at least two vertices and contributes deficiency at least \(\pi/2\) by the C3 auxiliary bound proved above. Their total already gives \(D\ge\pi\). If exactly one component is singular, total nullity at least two and the path/cycle recurrences force it to be a cycle of nullity two, of some length \(k\le N\) and rank \(k-2\).

The small cycles are settled as follows:

\begin{itemize}
\tightlist
\item
  \(k=3\): rank one forces all three lines to coincide, giving deficiency \(3\pi/2\).
\item
  \(k=4\): the two opposite orthogonal pairs are bases of one plane; the four edge deficiencies sum to \(\pi\).
\item
  \(k=5\): the five-cycle inequality proved above gives deficiency at least \(\pi\).
\end{itemize}

These cases include the \(N=4\) and \(N=5\) bases. If \(k\ge6\), apply the preceding cycle-projection argument. The resulting \(k-1\) lines lie in dimension at most \(k-3\) and \(k-1<N\). Embed their span into \(\mathbb R^{k-3}\) if necessary. The induction hypothesis gives projected deficiency at least \(\pi\), and projection did not increase deficiency. The original cycle therefore contributes at least \(\pi\) too. Every other component contributes a nonnegative amount.

This proves \(D\ge\pi\) for the sparse angle maximizer in all cases. Every other configuration has angle sum no greater, hence deficiency no smaller, proving the C5 bound for all configurations. Equality is attained by coordinate axes with multiplicities \((2,2,1,\ldots,1)\), which have exactly two coincident pairs.

\subsection{Provenance and verification}\label{provenance-and-verification}

The five proved statements above cover all their stated parameters and allow repeated lines. The proof is mathematical and does not rely on computation. For historical context, \href{https://www.math.u-szeged.hu/~vigvik/preprints/egyenesekszogei.pdf}{Fodor--V\'{i}gh--Zarn\'{o}cz} reports earlier small-dimensional results and identifies \href{https://doi.org/10.1007/BF02063286}{Fejes T\'{o}th's 1959 paper}; its full original proof was not inspected or used as a substitute for the arguments here. The C5 projection extension was reconstructed in this work; no novelty claim is made.

The companion Lean project checks selected analytic, geometric, and matrix lemmas, not the full general C1--C5 theorems. Its reproducibility files are \href{https://github.com/mpelteshki/climbing-to-the-frontier/blob/212bba53b1818c08c40b8c4b48f636651eaf6410/cagent/p1/verify.sh}{verification script}, \href{https://github.com/mpelteshki/climbing-to-the-frontier/blob/212bba53b1818c08c40b8c4b48f636651eaf6410/cagent/p1/evidence/replay/README.md}{fresh replay record}, and \href{https://github.com/mpelteshki/climbing-to-the-frontier/blob/212bba53b1818c08c40b8c4b48f636651eaf6410/cagent/p1/evidence/source-sha256.txt}{source checksums}. From the repository root, run \texttt{cd\ cagent/p1\ \&\&\ ./verify.sh}. C6 is excluded, and no organizer acceptance or platform submission is claimed.
\end{document}
